\documentclass[10pt,a4paper]{amsart}
\usepackage[margin=19mm]{geometry}
\usepackage{amsmath,amssymb,mathtools}
\usepackage[scr=rsfs,cal=euler]{mathalfa}
\usepackage{xfrac}
\usepackage[unicode,hidelinks,bookmarks=false]{hyperref}
\newcommand{\ie}{i.e.\ }
\newcommand{\cf}{cf.\ }
\newcommand{\resp}{respectively\ }
\newcommand{\Subsection}{\subsection}
\providecommand{\nobreakdash}{\nobreak\hbox{-}}
\setlength{\emergencystretch}{2em}
\multlinegap0pt
\usepackage{amssymb}
\newtheorem{lemma}{Lemma}
\title{A note on finite groups with few conjugacy classes of subgroups}
\author{Marius T\u{a}rn\u{a}uceanu}
\date{Source: arXiv:2607.05700v1 (CC BY 4.0)}
\begin{document}
\maketitle

\noindent\emph{Source and licence.} A note on finite groups with few conjugacy classes of subgroups. Version arXiv:2607.05700v1 (CC BY 4.0), distributed by arXiv under the Creative Commons Attribution 4.0 International licence, http://creativecommons.org/licenses/by/4.0/, as recorded in the arXiv licence metadata for this exact version and shown on the abstract page https://arxiv.org/abs/2607.05700v1. Full licence notice: \texttt{licenses/arxiv-2607.05700-CC-BY-4.0.txt}.\par\smallskip

\noindent\textit{Editorial scope.} Counted unit: the article's lemma lemma (source0.tex lines 143--145). The article's complete original proof of it is delivered with it (source0.tex lines 147--149, one \texttt{proof} environment of 3 lines). No mathematics is written, repaired or re-derived: the statement and the proof are the article's own lines, in the article's own notation, and no retained line is substituted or re-flowed.  No further source-local context is needed: every cross-reference of every retained line resolves inside the two delivered spans.  Nothing was omitted from any retained span, so the package carries no omission record.  No retained line contains a citation, so the wrapper carries no bibliography and no bibliography entry.  No retained line contains a figure environment, an image-inclusion command or drawing code, so no image is delivered and the package declares no asset dependency.  Editorial formatting only, in addition: an AMS \texttt{amsart} preamble that restates the article's own declarations and macros used by the retained lines, namely the environments \texttt{lemma} on separate counters, together with the standard packages those lines need.  The retained environments print in this document as Lemma 1 in order of appearance, whereas the article prints the same items with the numbering of its own full document.  The complete native source of the article as distributed in the arXiv e-print for this exact version is bundled unchanged as \texttt{sources/204563/source0.tex}.
\begin{lemma}
Let $p$ be a prime. Then $|{\rm PSL}(2,2^p)|$ is divisible by exactly three primes if and only if either $p=2$ or $p=3$.
\end{lemma}

\begin{proof}
We have $|{\rm PSL}(2,2^p)|=2^p(2^p-1)(2^p+1)$ and this number is divisible by exactly three primes if and only if both $2^p-1$ and $2^p+1$ are prime powers. This holds for $p=2$ and $p=3$. If $p\geq 5$, then $p$ is odd and so $2^p+1$ is divisible by $3$. On the other hand, in this case there exists a prime divisor $q$ of $2^p+1$ that does not divide $2^k+1$ for all $k\in\{1, \dots, p-1\}$ by Theorem B. Then $q\nmid 3$, that is $q\neq 3$, implying that $2^p+1$ cannot be a prime power. Thus the unique possibilities are $p=2$ and $p=3$.\qedhere
\end{proof}

\end{document}
